\section{Method}
\label{sec:method}

\subsection{Starting point: the MMAudio objective}
\label{sec:method:baseline}
Let $x_1 \in \R^{N \times d}$ be a VAE-encoded audio latent (for
the 16\,kHz small model, $N=250$, $d=20$). Conditioning consists of
CLIP tokens $z^{\text{clip}} \in \R^{64 \times 1024}$,
Synchformer tokens $z^{\text{sync}} \in \R^{192 \times 768}$, and
CLIP-text tokens $z^{\text{text}} \in \R^{77 \times 1024}$. With a
sampled time $t \sim \operatorname{LogNormal}$ and Gaussian
$x_0$, MMAudio minimizes the flow-matching loss of
Sec.\,\ref{sec:background},
\begin{equation}
  \LFM(\theta) = \E\,
    \bigl\lVert v_\theta(t, x_t; z^{\text{clip}}, z^{\text{sync}},
      z^{\text{text}}) - (x_1 - x_0) \bigr\rVert_2^2,
\end{equation}
with classifier-free dropout (probability 0.1) of each conditioning
stream. Gradients therefore flow through the MMDiT joint blocks
(joint QKV attention) and the Conditional Synchronization
Module (frame-aligned adaLN). Both mechanisms align video and
audio features \emph{pointwise}, as analyzed in
Sec.\,\ref{sec:background:mmaudio}. No other loss term is added
in \citet{cheng2024mmaudio}; the pointwise--relational gap we
target (Sec.\,\ref{sec:intro:pointwise}) lives precisely in the
space of objectives this expectation does not cover.

\subsection{The regularized objective}
\label{sec:method:reg}
We add a single relational term:
\begin{equation}
\label{eq:total-loss}
  \mathcal{L}(\theta) = \underbrace{\LFM(\theta)}_{\text{pointwise, MMAudio}}
    + \lambda_{\mathrm{GW}}(t_{\text{step}}) \cdot
    \underbrace{\LGW(\theta)}_{\text{relational, ours}},
\end{equation}
where $t_{\text{step}}$ is the global training iteration and
$\lambda_{\mathrm{GW}}(\cdot)$ is a schedule (Sec.\,\ref{sec:design}).
The two terms operate on orthogonal axes: $\LFM$ is a sum of
per-sample, per-time-step regression errors; $\LGW$ is a function
only of the batch-level $B \times B$ distance matrices and is
constant under any transformation that preserves pairwise
distances (translation, orthogonal rotation). A training step can
lower one without trivially lowering the other, which is what
makes the ablations in Sec.\,\ref{sec:experiments} informative.
The GW loss is computed \emph{per batch} from the two pairwise
distance matrices $\DV, \DA \in \R^{B \times B}$:
\begin{align}
  \DV[i,k] &= \lVert v_i - v_k \rVert_2^2, \\
  \DA[j,\ell] &= \lVert a_j - a_\ell \rVert_2^2,
\end{align}
where $v_i, a_j \in \R^{D}$ are the per-sample
\emph{representations} pulled from the network; the exact choice
of representation is design axis D1 (Sec.\,\ref{sec:design}).

\paragraph{Scale invariance.}
Because both $\DV$ and $\DA$ scale with the feature magnitude, and
because the ratio $\lVert \DV \rVert / \lVert \DA \rVert$ is
arbitrary, we normalize each by its mean before computing GW:
\begin{equation}
  \tilde D = D / \operatorname{mean}(D).
\end{equation}
This removes one confounder from the $\lambda_{\mathrm{GW}}$
sweep --- otherwise $\lambda_{\mathrm{GW}}$ would have to chase
feature-magnitude drift during training.

\subsection{Entropic GW via Peyr\'e mirror descent}
\label{sec:method:gw}
The relational loss is the entropic GW distance of
Eq.\,(\ref{eq:gw}) with uniform marginals $p = q = \mathbf{1}/B$:
\begin{equation}
  \LGW =
    \sum_{i,j,k,\ell}
    \bigl( \tilde\DV[i,k] - \tilde\DA[j,\ell] \bigr)^{2}
    T_{ij}\, T_{k\ell}
    - \varepsilon\, H(T),
\end{equation}
optimized over $T \in \Pi(p, q)$. We solve this with the
projected-gradient algorithm of \citet{peyre2016gromov}: at each
outer step, linearize the quadratic cost around the current $T$
using $G = -2\, \tilde\DV\, T\, \tilde\DA$, set
$K = \exp(-G / \varepsilon)$, and run $L$ Sinkhorn iterations
\citep{cuturi2013sinkhorn}:
\begin{equation}
  u \leftarrow p \oslash (K v),
  \qquad
  v \leftarrow q \oslash (K^{\top} u),
  \qquad
  T \leftarrow \operatorname{diag}(u)\, K\, \operatorname{diag}(v),
\end{equation}
with $\oslash$ element-wise division and a small floor
($10^{-8}$) to guard against division by zero. We use $K_{\text{out}}
= 5$ outer iterations and $L = 20$ inner Sinkhorn iterations by
default; these were fixed early and not retuned, since batches of
$B \le 512$ converge well before that.

Given the final coupling $\Tstar$, we return the full objective
$\sum_{ijk\ell}(\DV[i,k] - \DA[j,\ell])^2\, \Tstar_{ij}\,
\Tstar_{k\ell}$ evaluated as
$\lVert \tilde\DV \rVert^2_p + \lVert \tilde\DA \rVert^2_q
- 2 \langle \tilde\DV\, \Tstar\, \tilde\DA, \Tstar \rangle$, so
that the logged loss magnitude is comparable across variants.

\paragraph{Float32 is not optional.}
The inner Sinkhorn's ratios $p \oslash (K v)$ and the matrix
exponential $\exp(-G/\varepsilon)$ are both numerically fragile.
bfloat16 loses the dynamic range of $K$ as soon as $G$ gets large
(early training or small $\varepsilon$) and the Sinkhorn diverges.
We wrap the entire GW computation in
\texttt{torch.cuda.amp.autocast(enabled=False)} and cast inputs
to \texttt{float32} on entry. The rest of the network continues
to train in bfloat16. Empirically this is a measured cost: the
per-step overhead of GW is dominated by two $B \times B$
matrix--matrix multiplies per outer step, i.e. at $B{=}64$ it is
negligible relative to the transformer forward.

\subsection{Fused Gromov--Wasserstein variant}
\label{sec:method:fgw}
For the \texttt{fused} variant (D1), we additionally build a
pointwise cross-domain cost from $\ell_2$-normalized
representations,
$C_{\text{cross}}[i,j] = 1 - \hat v_i^{\top} \hat a_j$, and solve
the FGW problem \citep{vayer2020fused}:
\begin{equation}
  \LFGW(T) =
    \alpha \sum_{ijk\ell}
    (\tilde \DV[i,k] - \tilde \DA[j,\ell])^2 T_{ij} T_{k\ell}
    + (1-\alpha) \langle C_{\text{cross}}, T \rangle,
\end{equation}
with the same Peyr\'e--Sinkhorn machinery (the outer linearized
cost becomes $\alpha G_{\mathrm{GW}} + (1-\alpha) C_{\text{cross}}$).
We use $\alpha = 0.5$ by default.

\subsection{Real-video subset and DDP}
\label{sec:method:subset}
MMAudio's joint training interleaves audio--visual and
audio--text samples at approximately $1{:}1$
\citep{cheng2024mmaudio}; the MMAudio release we build on
preserves the same mix via a \texttt{MultiModalDataset} of
VGGSound (video + audio) plus audio-only corpora (AudioCaps,
AudioSetSL, BBCSound, FreeSound, Clotho). For audio-only
samples, the CLIP and Synchformer features are empty
placeholders, and the flow-matching stream feeds them as
$\varnothing_v, \varnothing_{\mathrm{syn}}$ tokens --- exactly
the upstream convention. GW, however, would compare empty-video
distances against real-audio distances, which is meaningless.
We therefore mask to the \texttt{video\_exist=True} subset
before building $\DV$, and skip the step entirely if the subset
has fewer than \texttt{min\_batch=4} samples. This is the only
place our loss treats the two data types differently from
upstream; flow matching continues to see the full batch.

When running under DDP, the GW loss is computed per-rank on the
local real-video subset and summed into the total loss before
\texttt{.backward()}. DDP's all-reduce then averages gradients
across ranks. This is an approximation --- the ``true''
batch-level GW on a global batch of size $B W$ would be a
different objective --- but it is cheap and consistent with the
rest of the training loop; a global-batch GW would require an
extra all-gather per step.

\subsection{What our term does that MMAudio's cannot}
\label{sec:method:what-changes}
Restating the positioning argument concretely, with the
motivating diagnostics of Sec.\,\ref{sec:intro:motivation} as
empirical context. All three
cross-modal mechanisms in MMAudio --- joint MMDiT attention, the
Conditional Synchronization Module, and empty-token masking ---
produce, at the level of the loss, terms of the form
$\mathcal{L}_i = \ell(f_\theta(z_i^{\text{video}}),
x_i^{\text{audio}})$, i.e. a function of a \emph{single} paired
$(z_i, x_i)$. In the decomposition of
\citet{wang2020understanding}, they are pure alignment terms.
$\LGW$ adds a term of the form
\[
  \mathcal{L}_{\mathrm{rel}} = \phi\bigl(
    \{ \lVert v_i - v_k \rVert \}_{i,k},
    \{ \lVert a_j - a_\ell \rVert \}_{j,\ell}
  \bigr),
\]
which is not expressible as a sum of per-pair terms. Two models
can have identical $\LFM$ and still differ arbitrarily in
$\LGW$: shuffle the per-sample offsets of the audio latent
manifold without changing paired distances, and $\LFM$ is
unchanged while $\LGW$ changes. This is the precise sense in
which the two loss terms are orthogonal.

\subsection{Engineering integration}
\label{sec:method:integration}
The GW loss is called from \texttt{Runner.train\_pass} in
\texttt{mmaudio/runner.py} \emph{after} the flow-matching forward,
and crucially \emph{outside} the \texttt{torch.compile}d
\texttt{train\_fn}. Compilation is a large speedup for the
transformer but would bake the Sinkhorn iterations into an inner
CUDA graph and make debugging numerics very painful; keeping GW
out of the compiled region costs little because its overhead is
dominated by the $B{\times}B$ matmuls, not kernel launch overhead.

Concretely, we snapshot the pre-CFG-null $\text{clip}_f$ before
\texttt{train\_fn} mutates it in place (the snapshot is needed
because the mutation fills rows with \texttt{empty\_clip\_feat}
for classifier-free guidance), normalize the audio latent with the
same $(\mu,\sigma)$ used inside the network, call
\texttt{compute\_gw\_regularization}, and add
$\lambda_{\mathrm{GW}}(t_{\text{step}}) \cdot \LGW$ to the total
loss. Effective $\lambda_{\mathrm{GW}}$ is logged alongside
$\LGW$ and $\LFM$ at every \texttt{log\_text\_interval}.
